% ============================================================
% paper/sections/07_continuity_k_measure.tex
% Continuity Measure k_EA (finalized)
% + FWHQ-01 foundational hardening additions
% ============================================================

\section{Continuity Measure $k_{\mathrm{EA}}$}

\subsection{Meaning of Continuity in OC}

Within the Ontology of Continua, continuity is not a geometric,
topological, or metric notion in the classical sense.
It denotes the degree of structural connectivity among admissible states
under the dynamics permitted by the continuum.

For the enterprise continuum, continuity expresses whether admissible
states remain mutually reachable without violating thresholds or
breaking required cycles.
Continuity therefore concerns definability of transitions, not quality
of outcomes.

Continuity is evaluated solely within the admissible region
$\Omega_{\mathrm{EA}}$.
Transitions that require leaving $\Omega_{\mathrm{EA}}$ are not
considered part of admissible connectivity.

\subsection{Definition of $k_{\mathrm{EA}}$}

The continuity measure $k_{\mathrm{EA}}$ is defined as a scalar quantity
bounded in the interval $[0,1]$.

A value of $k_{\mathrm{EA}} = 1$ indicates full structural connectivity
among admissible states under the defined dynamics.
A value of $k_{\mathrm{EA}} = 0$ constitutes a hard structural
discontinuity, at which point the enterprise continuum ceases to be a
valid object of diagnosis.

Intermediate values indicate partial connectivity.
They do not imply gradual degradation, recoverability, resilience, or
proximity to any desirable state.

The ordering of values in $(0,1)$ is not interpreted metrically.
Only boundary conditions and admissibility implications are relevant.

\subsection{Factorization of Continuity}

According to \texttt{ETS.K\_EA.v1.1}, the continuity measure
$k_{\mathrm{EA}}$ is factorized into a finite set of structural
components:

\begin{itemize}
  \item $H_{\Omega}$: indicator of non-empty admissible state space,
  \item $S_{\text{cycles}}$: contribution of required recurrent cycles,
  \item $S_{\text{obs}}$: contribution of admissible observability,
  \item $S_{\text{parent}}$: compatibility with parent continua,
  \item $S_{\text{thresholds}}$: compliance with admissibility
        thresholds.
\end{itemize}

The aggregation rule is multiplicative.
Violation of any hard component collapses the product and yields
$k_{\mathrm{EA}} = 0$.

This formulation enforces the principle that continuity is undefined if
any required structural condition fails.

\subsection{Hard-Zero Condition}

The factor $H_{\Omega}$ acts as a hard gate on continuity.
If the admissible state space becomes empty, continuity collapses
immediately, regardless of the values of all other components.

This reflects the principle that structural connectivity is undefined in
the absence of admissible states.
In such regimes, diagnostic statements have no formal meaning within the
model.

\subsection{Continuity Versus Performance}

Continuity must not be interpreted as a measure of performance,
efficiency, effectiveness, success, or health.

A system may exhibit high operational or financial performance while
having low structural continuity, and conversely.
No monotonic relationship between continuity and observed outcomes is
assumed or implied.

The sole function of $k_{\mathrm{EA}}$ is to determine whether diagnosis
remains structurally admissible within the model.

\subsection{Role in Phase Logic}

The value of $k_{\mathrm{EA}}$ influences phase logic indirectly through
its interaction with threshold components and termination conditions.

A decline of continuity toward zero signals approaching loss of
diagnosability.
However, $k_{\mathrm{EA}}$ does not prescribe interventions, targets, or
responses.

It serves exclusively as a structural indicator of admissibility within
the enterprise continuum, not as a control variable or optimization
objective.

%%%End of Document%%%
